\documentclass[11pt,a4paper]{article}
\usepackage[T1]{fontenc}
\usepackage[utf8]{inputenc}
\usepackage{lmodern,amsmath,amssymb}
\usepackage[margin=25mm]{geometry}
\usepackage{longtable,booktabs,array,calc}
\usepackage[hidelinks]{hyperref}
\hypersetup{pdftitle={A monotone hierarchy of upper bounds from one flowed correlator, and why no bound of this },pdfauthor={navstokgap research working notes}}
\newcommand{\tightlist}{\setlength{\itemsep}{0pt}\setlength{\parskip}{0pt}}
\setlength{\emergencystretch}{3em}
\setcounter{secnumdepth}{0}
\begin{document}
\hypertarget{a-monotone-hierarchy-of-upper-bounds-from-one-flowed-correlator-and-why-no-bound-of-this-kind-can-give-m0}{%
\section{\texorpdfstring{A monotone hierarchy of upper bounds from one
flowed correlator, and why no bound of this kind can give
\(m>0\)}{A monotone hierarchy of upper bounds from one flowed correlator, and why no bound of this kind can give m\textgreater0}}\label{a-monotone-hierarchy-of-upper-bounds-from-one-flowed-correlator-and-why-no-bound-of-this-kind-can-give-m0}}

Let \(\varphi_t\) be any flowed gauge-invariant local scalar with
\(\langle\varphi_t\rangle=0\) and let
\[C_0(\tau)=\int d^3z\;\big\langle\varphi_t(z,\tau)\,\varphi_t(0,0)\big\rangle^{\rm c}
=\int_{[m,\infty)}d\rho(E)\;e^{-E\tau/\hbar},\qquad \rho\ge0,\] be its
zero-spatial-momentum Euclidean correlator, with \(\rho\) the
Källén--Lehmann measure supplied by Osterwalder--Schrader
reconstruction. Writing \(M_k=\int E^k\,d\rho(E)\), every consecutive
ratio is an upper bound for the mass gap and the bounds improve
monotonically as the index decreases:
\[m\ \le\ \cdots\ \le\ \frac{M_{-1}}{M_{-2}}\ \le\ \frac{M_0}{M_{-1}}\ \le\ \frac{M_1}{M_0},
\qquad\text{with}\qquad \frac{M_1}{M_0}\ \text{the }f\text{-sum bound of}\]
\href{finiteness-half-flowed-susceptibility.md}{the finiteness note},
which is therefore the weakest member of the family. The first member
below it is
\[\boxed{\;m\ \le\ \frac{M_0}{M_{-1}}\ =\ \frac{\hbar\,C_0(0)}{\displaystyle\int_0^\infty C_0(\tau)\,d\tau}\;}\]
and it needs strictly less than the \(f\)-sum bound: no gradient term,
no equal-time Hamiltonian smearing, hence no hypothesis (F1). Both
quantities are integrals of the four-dimensional flowed correlator that
Lüscher's renormalization statement covers, so the only inputs are the
existence of the theory with the Osterwalder--Schrader axioms and
\(C_0\not\equiv0\). In free Maxwell theory the first three members take
the values \(1.504\), \(1.329\), \(1.128\) in units of
\(\hbar c/\sqrt t\), decreasing as the hierarchy requires. The negative
moments diverge when the spectral measure reaches down to \(E=0\), so
their finiteness is necessary for a gap; it is not sufficient, since
\(d\rho=e^{-\hbar c\kappa/E}dE\) has every negative moment finite with
\(\inf\operatorname{supp}\rho=0\). Consequently \textbf{no member of
this hierarchy, and no upper bound of this type, can establish \(m>0\)}:
the entire upper side of the problem settles the finiteness clause and
is silent on the clause that matters. Constants explicit; nothing
promoted.

\hypertarget{moments-of-the-flowed-correlator}{%
\subsection{1. Moments of the flowed
correlator}\label{moments-of-the-flowed-correlator}}

Assume the continuum theory exists with the Osterwalder--Schrader
axioms, a unique vacuum and a self-adjoint \(H\ge0\) with \(H\Omega=0\);
let \(\varphi_t\) be a gauge-invariant local scalar in the field at flow
time \(t>0\), with \(\langle\varphi_t\rangle_0=0\). Reconstruction
gives, for the zero-spatial-momentum correlator at Euclidean separation
\(\tau>0\),
\[C_0(\tau)=\int d^3z\,\big\langle\varphi_t(z,\tau)\varphi_t(0,0)\big\rangle^{\rm c}_0
=\big\langle\Omega,\Phi^\dagger e^{-\tau H/\hbar}\Phi\,\Omega\big\rangle
=\int_{[m,\infty)}e^{-E\tau/\hbar}\,d\rho(E),\] where \(\Phi\) is the
zero-momentum smeared operator of
\href{finiteness-half-flowed-susceptibility.md}{the finiteness note} §2
and \(d\rho\) is the spectral measure of \(H\) in the state
\(\Phi\Omega\), divided by the volume so that \(\rho\) is intensive. The
vacuum contribution is removed by the connected part, so
\(\operatorname{supp}\rho\subset[m,\infty)\) where
\(m=\inf\big(\operatorname{spec}H\setminus\{0\}\big)\). Set
\[M_k=\int_{[m,\infty)}E^k\,d\rho(E)\in[0,\infty],\qquad k\in\mathbb Z .\]
Two of them have elementary readings: \[M_0=C_0(0^+),\qquad
M_{-1}=\frac1\hbar\int_0^\infty C_0(\tau)\,d\tau,\qquad
M_{1}=-\hbar\,\frac{dC_0}{d\tau}\Big|_{\tau=0^+},\] the equal-time
zero-momentum variance, the time-integrated correlator, and the
\(f\)-sum rule.

\hypertarget{the-hierarchy}{%
\subsection{2. The hierarchy}\label{the-hierarchy}}

\textbf{Proposition 1.} If \(M_k\) and \(M_{k+1}\) are finite and
\(M_k>0\), then \[m\ \le\ \frac{M_{k+1}}{M_k}.\] If in addition
\(M_{k-1}<\infty\) and \(M_{k-1}>0\), then
\[\frac{M_k}{M_{k-1}}\ \le\ \frac{M_{k+1}}{M_k}.\]

\emph{Proof.} On \(\operatorname{supp}\rho\) one has \(E\ge m\), so
\(M_{k+1}=\int E\cdot E^k\,d\rho\ge m\int E^k\,d\rho=mM_k\), which is
the first claim. For the second, Cauchy--Schwarz in \(L^2(d\rho)\)
applied to \(E^{(k-1)/2}\) and \(E^{(k+1)/2}\) gives
\(M_k^2=\big(\int E^{(k-1)/2}E^{(k+1)/2}d\rho\big)^2\le M_{k-1}M_{k+1}\),
that is, \((M_k)\) is log-convex; dividing by \(M_{k-1}M_k\) gives the
claim. \(\square\)

The chain therefore runs downward from the \(f\)-sum bound:
\[m\ \le\ \cdots\ \le\ \frac{M_{-1}}{M_{-2}}\ \le\ \frac{M_{0}}{M_{-1}}\ \le\ \frac{M_{1}}{M_{0}} .\]
Each step uses one more inverse power of the energy and so weights the
low-lying states more heavily; in the limit of small index the ratio
converges to \(\inf\operatorname{supp}\rho\) whenever the moments stay
finite, by dominated convergence applied to \(\rho\) restricted to
\([m,m+\epsilon]\) against its complement.

\hypertarget{the-first-improvement-and-what-it-costs}{%
\subsection{3. The first improvement, and what it
costs}\label{the-first-improvement-and-what-it-costs}}

\textbf{Corollary 2.} If \(0<M_{-1}\) and \(M_0<\infty\), then
\[m\ \le\ \frac{M_0}{M_{-1}}=\frac{\hbar\,C_0(0^+)}{\int_0^\infty C_0(\tau)\,d\tau}\ <\ \infty .\]

Compare the inputs.

\begin{longtable}[]{@{}
  >{\raggedright\arraybackslash}p{(\columnwidth - 4\tabcolsep) * \real{0.3333}}
  >{\raggedright\arraybackslash}p{(\columnwidth - 4\tabcolsep) * \real{0.3333}}
  >{\raggedright\arraybackslash}p{(\columnwidth - 4\tabcolsep) * \real{0.3333}}@{}}
\toprule\noalign{}
\begin{minipage}[b]{\linewidth}\raggedright
bound
\end{minipage} & \begin{minipage}[b]{\linewidth}\raggedright
needs
\end{minipage} & \begin{minipage}[b]{\linewidth}\raggedright
supplied by
\end{minipage} \\
\midrule\noalign{}
\endhead
\bottomrule\noalign{}
\endlastfoot
\(M_1/M_0\) & \(\int\langle|\delta O/\delta A|^2\rangle<\infty\),
equal-time spatial smearing & hypothesis (F1) of the finiteness note \\
\(M_0/M_{-1}\) & \(C_0(0^+)<\infty\) and \(\int_0^\infty C_0>0\) &
four-dimensional flow: Lüscher's statement \\
\end{longtable}

The \(f\)-sum bound required the Hamiltonian counterpart of the flow's
renormalization, because \(M_1\) involves the functional derivative of
the operator with respect to the equal-time field. The ratio
\(M_0/M_{-1}\) involves only the flowed correlator itself, and Lüscher
states that the gauge field at flow time \(t>0\) is a smooth
renormalized field, so that expectation values of local gauge-invariant
expressions in it are well-defined physical quantities
(arXiv:1006.4518v3, abstract and §2; passage level via the
\href{../docs/Luscher_WilsonFlow_1006.4518v3.md}{local companion}).
Hypothesis (F1) is therefore not needed for Corollary 2, and the
remaining inputs are the existence of the theory with the axioms, which
is T4 of \href{mass-gap-obligations-lattice.md}{the obligations map},
and \(C_0\not\equiv0\).

\textbf{Proposition 3 (the finiteness half, restated).} Assume T4 and
that some flowed local gauge-invariant scalar has \(C_0\not\equiv0\)
with \(C_0(0^+)<\infty\). Then \(m<\infty\).

\emph{Proof.} \(C_0\not\equiv0\) and \(\rho\ge0\) give \(M_{-1}>0\);
Corollary 2 applies. \(\square\)

The clause \(m<\infty\) of the Jaffe--Witten statement is thus
equivalent, given the construction, to the nontriviality of a single
flowed correlator, which is the sharpest form this programme has
reached.

\hypertarget{free-field-check-of-the-monotonicity}{%
\subsection{4. Free-field check of the
monotonicity}\label{free-field-check-of-the-monotonicity}}

For free Maxwell theory with \(\varphi_t\) the flowed magnetic energy
density, \href{flowed-bound-free-field.md}{the free-field note} gives
the spectral weight \(d\rho\propto k^4e^{-4tk^2}dk\) at energy
\(E=2\hbar ck\). With \(\int_0^\infty k^ne^{-4tk^2}dk\) evaluated by
\(\int_0^\infty k^{2j}e^{-\alpha k^2}dk=\frac{(2j-1)!!}{2(2\alpha)^j}\sqrt{\pi/\alpha}\)
and \(\int_0^\infty k^{2j+1}e^{-\alpha k^2}dk=\frac{j!}{2\alpha^{j+1}}\)
at \(\alpha=4t\):
\[\frac{M_1}{M_0}=2\hbar c\,\frac{\int k^5e^{-4tk^2}}{\int k^4e^{-4tk^2}}
=\frac{8}{3\sqrt\pi}\,\frac{\hbar c}{\sqrt t}\approx1.504\,\frac{\hbar c}{\sqrt t},\]
\[\frac{M_0}{M_{-1}}=2\hbar c\,\frac{\int k^4e^{-4tk^2}}{\int k^3e^{-4tk^2}}
=\frac{3\sqrt\pi}{4}\,\frac{\hbar c}{\sqrt t}\approx1.329\,\frac{\hbar c}{\sqrt t},\]
\[\frac{M_{-1}}{M_{-2}}=2\hbar c\,\frac{\int k^3e^{-4tk^2}}{\int k^2e^{-4tk^2}}
=\frac{2}{\sqrt\pi}\,\frac{\hbar c}{\sqrt t}\approx1.128\,\frac{\hbar c}{\sqrt t},\]
using \(\int k^5=1/(64t^3)\), \(\int k^4=3\sqrt\pi/(256\,t^{5/2})\),
\(\int k^3=1/(32t^2)\), \(\int k^2=\sqrt\pi/(32\,t^{3/2})\). The three
numbers decrease, as Proposition 1 requires, and the first reproduces
the free-field note. In this theory \(M_k\) is finite for \(k\ge-4\) and
diverges for \(k\le-5\), since \(d\rho\propto E^4dE\) near \(E=0\): the
hierarchy terminates, and its terminating member does not reach \(m=0\).
Termination at a finite index is the spectral signature of weight
accumulating at zero energy.

\hypertarget{why-the-upper-side-cannot-give-m0}{%
\subsection{\texorpdfstring{5. Why the upper side cannot give
\(m>0\)}{5. Why the upper side cannot give m\textgreater0}}\label{why-the-upper-side-cannot-give-m0}}

Every bound above is of the form \(m\le(\text{ratio of moments})\), and
each is finite whenever the corresponding moments are. Finiteness of all
negative moments is necessary for \(m>0\) but not sufficient: take
\[d\rho(E)=e^{-\hbar c\kappa/E}\,dE\ \ \text{near}\ E=0,\qquad\kappa>0,\]
whose support reaches \(E=0\), so \(m=0\), while
\(M_{-k}=\int E^{-k}e^{-\hbar c\kappa/E}dE=(\hbar c\kappa)^{1-k}\Gamma(k-1)\)
after \(u=\hbar c\kappa/E\), finite for every \(k\ge2\). A theory with
this spectral density would satisfy every bound in the hierarchy with a
positive right-hand side and still be gapless. No amount of information
about the moments of one correlator distinguishes it from a gapped
theory.

The variational principle produces upper bounds, and \(m>0\) is a lower
bound; the two require different mathematics. The entire content of
Sections 1--4, like that of \href{lattice-gap-upper-bounds.md}{the
upper-bound note} and \href{polyakov-average-gap-bound.md}{the Polyakov
note}, belongs to the finiteness clause. The clause \(m>0\) needs an
argument of the kind attempted on the lower side: an operator inequality
\(H\ge m\,(1-|\Omega\rangle\langle\Omega|)\), which is what the
strong-coupling theorem of \href{strong-coupling-uniform-gap.md}{the T2
note} achieves at fixed cutoff and what
\href{schur-error-ultraviolet.md}{the Schur note} shows is obstructed by
ultraviolet dressing in the continuum.

\hypertarget{consequence-for-state}{%
\subsection{6. Consequence for STATE}\label{consequence-for-state}}

The upper side is closed as far as it can go: a monotone hierarchy whose
best members need only T4 and the nontriviality of one flowed
correlator, with the \(f\)-sum bound identified as its weakest member
and hypothesis (F1) eliminated. The finiteness clause is therefore
reduced to the construction plus one nonvanishing correlator, and no
further work on upper bounds changes the standing of the problem.
Everything remaining is the lower bound: T2\('\) at fixed cutoff and its
uniformity, which is where the next steps belong.
\end{document}
